% =====================================================================
%  Night Flash Edge -- Technical Report
%  Build:  ./build.sh          (or: tectonic -X compile technical-report.tex)
% =====================================================================
\documentclass[11pt,a4paper]{article}

\usepackage[margin=25mm]{geometry}
\usepackage{amsmath,amssymb,booktabs,graphicx,longtable,array,enumitem}
\usepackage[hidelinks]{hyperref}
\usepackage{fancyhdr,titlesec,microtype}
\usepackage{xcolor}
\usepackage{listings}

\definecolor{codebg}{RGB}{247,247,249}
\definecolor{rulegray}{RGB}{200,200,205}
\definecolor{accent}{RGB}{30,60,110}

\lstset{
  backgroundcolor=\color{codebg},
  basicstyle=\ttfamily\small,
  breaklines=true,
  frame=single,
  rulecolor=\color{rulegray},
  columns=fullflexible,
  keepspaces=true,
  showstringspaces=false,
}
\lstdefinelanguage{nfconsole}{
  morekeywords={nf,python,pip},
  sensitive=true,
  morecomment=[l]{\#},
}

\titleformat{\section}{\Large\bfseries\color{accent}}{\thesection}{0.7em}{}
\titleformat{\subsection}{\normalsize\bfseries}{\thesubsection}{0.6em}{}
\titlespacing*{\section}{0pt}{14pt}{6pt}
\titlespacing*{\subsection}{0pt}{10pt}{4pt}

\pagestyle{fancy}
\fancyhf{}
\fancyhead[L]{\small\color{gray}Night Flash Edge --- Technical Report}
\fancyhead[R]{\small\color{gray}\thepage}
\renewcommand{\headrulewidth}{0.4pt}

\newcommand{\code}[1]{\texttt{#1}}
\newcommand{\flag}[1]{\texttt{--#1}}

\title{\vspace{-2.5cm}\textbf{Night Flash Edge}\\[6pt]
\large A Compact Unpaired Style-Transfer System for a Disposable-Camera Aesthetic\\[4pt]
\large Technical Report}
\author{}
\date{\today}

\begin{document}
\maketitle
\vspace{-1.2cm}

\begin{abstract}
\noindent
This report documents a compact learned image-translation model intended to
reproduce a disposable-camera / Night Flash aesthetic inside an existing native
camera application. The system is deliberately small enough to run in real time on
a phone and to export to Core ML, and it learns from \emph{unpaired} image
domains, so no aligned before/after photographs are required.

The report covers the architecture and its three switchable capabilities; the
unpaired training objective and a correction to it that materially changed
training behaviour; a synthetic degradation pipeline that manufactures a source
domain from the target photographs themselves, in two variants with different
guarantees; a held-out validation regime with reproducible strips, a metric
history and best-checkpoint selection; the inference path including optional
tiled and downscaled modes; and the measurement results behind each of these
decisions. Measurements quoted throughout were produced by the code in this
repository and are reproducible with the scripts referenced in
\S\ref{sec:reproduce}.
\end{abstract}

\tableofcontents
\newpage

% =====================================================================
\section{Problem statement and constraints}

\subsection{What the model must do}

Given an ordinary photograph, produce one that reads as a disposable-camera or
Night Flash capture: hard direct flash on a near subject, rapid luminance falloff
to a dark surround, crushed shadows, warm ambient tint, visible film grain and
light leak. Two properties are non-negotiable:

\begin{enumerate}[leftmargin=1.4em]
  \item \textbf{Geometry, faces and text must survive.} The model is a filter,
        not a generator. A face that changes shape is a defect, not a style.
  \item \textbf{It must run on a phone, in real time.} That caps the parameter
        budget at single-digit millions and forbids any heavyweight auxiliary
        network in the shipped graph.
\end{enumerate}

\subsection{Why unpaired}

Genuine before/after pairs of the same scene do not exist. The style lives in
lighting and film response, not in content, so the two domains share
\emph{content} statistics only loosely. This is unpaired domain translation in
the CycleGAN sense, and it imposes a structural constraint on the loss: nothing
in the objective may compare a generated image to a specific target pixel.

\subsection{Deployment boundary}

The application already owns the Vision person mask and the Core Image / Metal
finishing chain. The learned model is inserted before that chain and is
responsible for the learned appearance only. Grain, leak and vignette are
rendered by a seeded procedural stage that can be disabled at export time, which
keeps the deterministic and stochastic paths independently testable.

% =====================================================================
\section{Architecture}

\subsection{Generator}

The default architecture, \code{edgehat}, is a convolutional encoder/decoder
with two kinds of attention.

\begin{description}[leftmargin=1.4em,style=nextline]
  \item[\code{EdgeViT} blocks] Multi-head self-attention confined to local
        windows. Linear in image size, and the window can be precomputed at
        construction time so nothing depends on a traced tensor shape.
  \item[HAT block] One cross-attention stage with overlapping shifted windows,
        giving a global receptive field for the price of a single block.
\end{description}

Level widths are $48 \rightarrow 96 \rightarrow 192 \rightarrow 384 \rightarrow 384$,
with a single fusion at the third stage. The total is \textbf{4.02 M parameters}
for the \code{freqsplit} head.

\begin{table}[h]
\centering
\begin{tabular}{@{}llrl@{}}
\toprule
Architecture & Attention & Params & Role \\
\midrule
\code{edgehat} & local windows $+$ one HAT block & 4.02 M & default \\
\code{swinir}  & windowed, relative position bias & 4.92 M & comparison arm \\
\code{restormer} & over channels & 0.76 M & legacy \\
\code{hybrid}  & pooled tokens & 0.76 M & legacy, keeps old checkpoints \\
\bottomrule
\end{tabular}
\caption{Architectures. The two legacy entries do not honour the newer heads and
exist only so pre-existing checkpoints remain loadable.}
\label{tab:archs}
\end{table}

\subsection{Bounded residual, and why it is the load-bearing decision}

The network predicts a residual, and the output is
\begin{equation}
  x_{\text{out}} = \operatorname{clamp}\!\left(x_{\text{in}} + \Delta(x),\, 0,\, 1\right),
\end{equation}
with $\lVert\Delta\rVert_\infty$ capped by \code{residual\_limit}.

This single constraint is what protects identity. The model is structurally
incapable of inventing content, because the only way to produce a pixel is to
perturb the pixel that was already there. It also means an under-trained model
degrades to a near-identity filter rather than to noise, which is the correct
failure direction for a camera filter.

\subsection{Output heads}

The head decides what the network is permitted to do.

\begin{table}[h]
\centering
\small
\begin{tabular}{@{}llrl@{}}
\toprule
Head & Learns & Params & Character \\
\midrule
\code{tanh}       & bounded delta                      & 21 k  & safest with little data \\
\code{perpixel}   & per-pixel delta, clamped on read   & 42 k  & free to repaint pixels \\
\code{freqsplit}  & low / mid / high bands separately  & 63 k  & grain and halation in the high band \\
\code{lut3d}      & $33^3$ trilinear LUT $+$ per-pixel & 129 k & global grade and local detail \\
\midrule
\code{lut3d+freqsplit}  & LUT and band split           & 192 k & recommended \\
\code{lut3d+perpixel}   & LUT and per-pixel            & 171 k & \\
\code{tanh+freqsplit}   & bounded and band split       & 84 k  & everything bounded \\
\bottomrule
\end{tabular}
\normalsize
\caption{Heads and composites. \code{perpixel+freqsplit} is deliberately absent:
both are per-pixel, so it costs parameters without adding a distinct inductive
bias.}
\label{tab:heads}
\end{table}

The default is \code{freqsplit} rather than a pure per-pixel head because grain
and halation live in the high frequency band, which a per-pixel parameterisation
represents poorly. Splitting the residual by frequency lets each band be
predicted at the resolution it needs.

\subsection{Three capabilities}

Each is independently switchable and each is gated by a \emph{zero-initialised}
mix scalar, so a freshly built model is exactly the plain residual and the extras
can only engage once training gives them a reason to. That is what makes them safe
to enable by default.

\begin{description}[leftmargin=1.4em,style=nextline]
  \item[Halation] A physical operator, not a learned guess:
    \begin{equation}
      \begin{aligned}
        \ell      &= Y \cdot (0.2126,\, 0.7152,\, 0.0722) \\
        h        &= \operatorname{relu}(\ell - \tau)\,/\,(1-\tau) \\
        s        &= \operatorname{DilateBlur}(h) \\
        g        &= \operatorname{relu}(s + T(F)) \odot \mathbf{tint}
      \end{aligned}
    \end{equation}
    Light scattering through the film base is a low-pass operation on
    over-threshold luminance, so a dilated depthwise blur is both faithful and
    very cheap. A learned convolution stack would be strictly worse: more
    parameters, less physics, and no ability to generalise to a highlight the
    training set never contained.
  \item[Subject gate] A soft $S(x,y) \in [0,1]$ from mid-level features. When a
    Vision person mask is supplied it is adaptively pooled to the same grid and
    added, so the gate is \emph{guided} rather than guessed. Tone is then
    modulated as $\Delta + \lambda S \tanh(\Delta)$, letting the subject diverge
    from the background without either running away.
  \item[Region exposure] Pools to a $4\times4$ grid and applies a $1\times1$
    convolution, so each cell is conditioned on its own features. Flattening the
    grid into one linear layer would couple every cell to every other.
\end{description}

\paragraph{Interaction with tiled inference.}
Region exposure conditions each cell on that cell's own features, which assumes
the grid spans the photograph. Under tiling it spans a tile. \code{Session.infer}
prints a note when it detects tiling together with region exposure; either infer
whole-frame or train with \flag{no-region-exposure}.

% =====================================================================
\section{The training objective}

\subsection{Unpaired baseline}

The generator is trained against a PatchGAN discriminator with a set of
constraints:

\begin{equation}
  \mathcal{L}_G = \mathcal{L}_{\mathrm{adv}}
    + 1.5\,\mathcal{L}_{\mathrm{content}}
    + 0.35\,\mathcal{L}_{\mathrm{moment}}
    + 0.75\,\mathcal{L}_{\mathrm{identity}}
    + 0.02\,\mathcal{L}_{\mathrm{effect}} .
  \label{eq:baseline}
\end{equation}

The discriminator is the only term that sees targets as targets. Everything else
is a leash. These weights are fixed and are not exposed on \code{nf} or
\code{Session}; the legacy \code{nightflash-train} has flags for them.

\subsection{A correction: where the texture penalty belonged}

The content term is a convex combination of an $L_1$ half and a Sobel edge half.
In the original implementation both halves were evaluated on the
\emph{post-effect} image \emph{$x_{\text{styled}}$}, the output after grain and
halation had been applied.

That is the wrong place for the $L_1$ half, and measurably so. $L_1$ is texture
sensitive, so comparing a grainy image against a clean source charges the output
for the grain. Measured on a trained model:

\begin{center}
\begin{tabular}{@{}lr@{}}
\toprule
Quantity & Value \\
\midrule
$\mathrm{content}\!\left(x_{\text{det}}, x_{\text{src}}\right)$ & 0.01451 \\
$\mathrm{content}\!\left(x_{\text{styled}}, x_{\text{src}}\right)$ & 0.04237 \\
\textbf{Penalty the stochastic effects incur} & \textbf{$+0.02786$} \\
\midrule
$\lVert x_{\text{det}} - x_{\text{src}}\rVert_1$ & 0.00891 \\
$\lVert x_{\text{styled}} - x_{\text{det}}\rVert_1$ & 0.01601 \\
\bottomrule
\end{tabular}
\end{center}

The effects were being charged roughly \emph{1.9 times} more than the entire
learned residual was charged for performing the whole style transfer. With the
adversarial term at weight $1.0$ asking for grain and the content term at weight
$1.5$ penalising it, the content term simply won, and the model learned to
suppress the grain and halation the effect head exists to produce. This is the
classical failure where an $L_1$ or $L_2$ penalty on a stochastic target averages
the stochastic component out of both sides.

\paragraph{The fix.} Split the two halves across different images:
\begin{equation}
  \mathcal{L}_{\mathrm{content}}
    = 0.65 \cdot L_1\!\left(x_{\text{det}},\, x_{\text{src}}\right)
    + 0.35 \cdot \mathcal{E}\!\left(x_{\text{styled}},\, x_{\text{src}}\right).
  \label{eq:content}
\end{equation}

The $L_1$ half now compares the \emph{pre-effect} image, so stochastic grain is
not charged. The edge half still compares the \emph{styled} image, because
structure genuinely should be verified on the final output, and a Sobel response
is indifferent to fine texture while being sensitive to broken geometry.

\paragraph{Consequence.} Grain and halation become visibly stronger, which is the
intent. It also means checkpoints trained under the old balance are not comparable
to new ones, so \code{LOSS\_VERSION} was incremented to $2$ and is written into
\code{config.json} and every checkpoint. Resuming a checkpoint whose recorded
version differs emits a warning rather than silently mixing two objectives.

\subsection{Supervised reconstruction, where a correspondence exists}

When the training mode supplies a known corresponding target, one more term is
available:
\begin{equation}
  \mathcal{L}_{\mathrm{paired}}
    = L_1\!\left(x_{\text{det}},\, y\right)
    + 0.5 \cdot \mathcal{E}\!\left(x_{\text{det}},\, y\right),
\end{equation}
evaluated on the \emph{deterministic} output. Pointing an $L_1$ term at a
stochastic target would re-create exactly the problem
Eq.~\eqref{eq:content} exists to avoid, so it is not done.

% =====================================================================
\section{Synthetic degradation}
\label{sec:degrade}

\subsection{The problem it solves}

Unpaired translation needs a source domain. When only a target domain is
available, the source has to be manufactured. This module simulates what the
same scene would have looked like through a non-flash camera.

\subsection{Pipeline and ordering}

The stages run in the order a camera pipeline would be \emph{undone}. Resampling
happens first so that photometric operations act on clean pixels.

\begin{center}
\begin{tabular}{@{}cl@{}}
\toprule
\# & Stage \\
\midrule
1 & crop and small rotation (shared with the target) \\
2 & \textbf{downscale then upscale} --- removes grain and fine detail \\
3 & Gaussian blur --- residual softening \\
4 & contrast compression about mid-grey \\
5 & black lift, $x \mapsto x + b(1-x)$, so white stays white \\
6 & exposure gain \\
7 & white-balance shift, warmed \\
8 & saturation increase \\
9 & light sensor noise \\
\bottomrule
\end{tabular}
\end{center}

Stages 4 and 5 carry most of the domain shift: compressing the tonal range is
what flattens the flash falloff signature, and it is the stage that makes the
source look like an evenly-lit non-flash capture rather than a dark one.

\subsection{Measured effect}

On a real photograph, averaged over eight draws per profile:

\begin{center}
\begin{tabular}{@{}lccc@{}}
\toprule
 & \code{subtle} & \code{moderate} & \code{strong} \\
\midrule
high-frequency energy (relative) & 0.52 & 0.30 & 0.21 \\
tonal range $p_{99}-p_1$         & 0.75 & 0.57 & 0.49 \\
saturation (absolute)            & 0.04 & 0.08 & 0.12 \\
warmth $R-B$ (absolute)          & 0.040 & 0.078 & 0.121 \\
\bottomrule
\end{tabular}
\end{center}

\subsection{Two honest limitations}

\paragraph{Resampling cannot remove a wide halo.}
An early draft of this document claimed that resampling destroys halation bloom.
Measurement refutes it. Halation is a wide, \emph{low-frequency} phenomenon, and a
low-pass filter has no problem to solve: blurring a glow makes it brighter by
spreading the core into the ring. On a synthetic probe with a specular point and a
halo, the ratio of ring luminance to peak luminance \emph{rises} from $0.57$ to
$0.86$ as blur goes from $0$ to $4$. What the pipeline achieves is attenuation,
driven by the tonal compression rather than the filtering, and that is sufficient
for the purpose: the model sees a highlight with less glow than the target has and
learns the difference. It would not be sufficient if the source needed to contain
\emph{zero} halation, which would require re-rendering the lighting rather than
filtering it.

\paragraph{The source must remain plausible.}
A first \code{strong} preset used blur up to $2.5$ and black lift up to $0.20$.
The result looked foggy rather than like a photograph, and its added noise was
glaring against the softened background. Training on that teaches a model to
expect an input distribution that never occurs at inference, which is worse than
a weaker domain shift. The ranges were pulled back until \code{strong} still reads
as a believable overcast capture.

\subsection{Variation}

Three independent levels, plus two optional mechanisms:

\begin{enumerate}[leftmargin=1.4em]
  \item \textbf{Per image} --- every parameter drawn independently per item.
  \item \textbf{Per epoch} --- free. The dataset seeds its generator with the
        epoch index, so the draw differs every epoch with no extra machinery.
  \item \textbf{Per run} --- the run seed shifts the whole stream.
  \item \textbf{Identity path}, probability \code{none\_prob} (default $0.10$):
        the target is fed through undegraded. This anchors the identity case and
        widens the input distribution, which is the main defence against the
        failure mode in \S\ref{sec:inverse}.
  \item \textbf{Curriculum} (on by default): the profile ramps
        \code{subtle} $\rightarrow$ \code{moderate} $\rightarrow$ \code{strong}
        across training. Early epochs get an easy mapping, later epochs get
        robustness.
\end{enumerate}

% =====================================================================
\section{Training modes}

\subsection{Three ways to form a pair}

\begin{center}
\begin{tabular}{@{}llll@{}}
\toprule
\code{train\_mode} & Source & Target & Correspondence \\
\midrule
\code{unpaired} & \code{source\_dir} & \code{target\_dir} & none, cross-domain \\
\code{paired}   & $\mathrm{deg}(x_i)$ & $x_i$ & exact, same file \\
\code{sameset}  & $\mathrm{deg}(x_j)$, $j \neq i$ & $x_i$ & none, same pool \\
\bottomrule
\end{tabular}
\end{center}

In both degradation modes the \emph{shared geometric warp is applied first, to
both images}, and the degradation runs afterwards on the source only. That
ordering is what keeps the two halves consistent: any crop, flip or rotation is
shared, and only the simulated camera differs.

\subsection{\texorpdfstring{\code{sameset}}{sameset} removes a shortcut}
\label{sec:inverse}

In \code{paired} mode a shortcut is available to the network: learn to invert
\emph{this specific} corruption. The cheapest way to reduce the loss is then
$\mathrm{deg}^{-1}$, and the model never learns anything about flash physics that
generalises. It will look excellent on freshly-degraded images and poor on a real
photograph, and the failure is invisible unless you look for it.

In \code{sameset} that shortcut does not exist. $x_i$ and $x_j$ are unrelated
photographs, so there is no pointwise relationship to invert and the model must
learn the conditional distribution of the look instead.

\paragraph{The tell.} If the model produces good results on freshly degraded
images but poor results on ordinary photographs, it has learned the inverse rather
than the aesthetic. The validation strip of \S\ref{sec:val} makes this visible
within a few epochs.

\paragraph{The cost.} Without a correspondence, \code{sameset} cannot report a
real quality number. Only distribution-level metrics apply, and \code{best.pt} is
deliberately not written, because ranking checkpoints on a PSNR that compares two
unrelated photographs would be meaningless.

\subsection{Mode resolution}

\code{train\_mode} defaults to \code{None} and is resolved from the arguments
actually supplied: a single pool resolves to \code{paired}, two distinct pools
resolve to \code{unpaired} with a printed notice. An explicitly requested mode is
always honoured and rejected loudly when it cannot work. This keeps
\code{paired} as the default for new single-pool work without breaking existing
two-pool scripts.

\subsection{A shared-warp bug worth recording}

The first implementation of \code{make\_pair} sampled scale, rotation and flips
once and then applied them to the source and the target in two separate calls.
The crop \emph{offset}, however, was drawn from the shared random stream inside
each call, so the two images received different offsets. This is harmless in
\code{unpaired} mode, where the images are unrelated anyway. In \code{paired} mode
it destroys the entire point: the correspondence.

The offset is now sampled once as a \emph{fraction} and clamped independently per
image, so two differently sized images still receive the same relative crop.
Verified: an identical input image now yields byte-identical source and target in
60 of 60 trials.

% =====================================================================
\section{Held-out validation}
\label{sec:val}

\subsection{Why a strip and not a number}

A loss curve shows the objective is decreasing. It does not show whether the
output still looks like a photograph. One strip image per run, rendered every
epoch, is cheap and is the thing a person actually looks at.

\subsection{Layout}

One image per row, three cells per image, ordered \code{input} $\vert$
\code{output} $\vert$ \code{target}. Rows rather than columns so a given image's
input and output are adjacent and the comparison needs no eye travel. The header
carries the epoch, the mode, the active degradation profile and the aggregate
metrics.

\subsection{Two determinisms, and why both are required}

\begin{description}[leftmargin=1.4em,style=nextline]
  \item[\emph{Which} images appear.] \code{val\_select="fixed"} (the default)
        keeps the same set every epoch so per-image progress is readable.
        \code{"random"} answers a different question --- is the model still healthy
        on \emph{arbitrary} images --- and hides the failure where a model looks
        great on a handful it has memorised.
  \item[\emph{How} each is degraded and seeded.] Always pinned, unconditionally.
        Seeds derive from the filename via SHA-256, not from the
        epoch-dependent stream, so a strip rendered in one session matches one
        rendered last week.
\end{description}

The second is not optional. A validation image degraded with a fresh random draw
each epoch is a different problem each time, and any apparent improvement is
noise. Python's \code{hash()} is salted per interpreter and therefore unusable
here; a cryptographic digest is required for cross-session stability.

\subsection{Splitting}

Validation images are drawn from the target pool by a deterministic, seeded
split and are never trained on. An explicit \flag{val-holdout-dir} overrides the
fraction. A pool too small to split leaves the validation half empty rather than
silently stealing training data.

\subsection{Metric history and checkpoint selection}

One CSV row per epoch:
\code{epoch, images, psnr, ssim, l1, effect}. This serves two purposes. It plots
the curve, and more importantly it makes checkpoint selection possible at all:
before this, only \code{latest.pt} existed, so a run that peaked at epoch $40$ and
decayed by epoch $60$ had no way to recover the good weights.

\begin{itemize}[leftmargin=1.4em]
  \item \code{best.pt} is written whenever validation PSNR improves, in
        \code{paired} mode only, for the reason in \S\ref{sec:inverse}.
  \item \code{keep\_last} prunes older per-epoch checkpoints. Sixty epochs of a
        $17$ MB checkpoint is a gigabyte, which fills Kaggle's disk and makes
        every snapshot slow.
\end{itemize}

\subsection{Exponential moving average}

\code{ema\_decay} maintains a shadow copy of the generator weights and is
checkpointed and resumed with the run. Averaging over the trajectory is standard
for adversarial training and reduces frame-to-frame flicker, which matters
directly for a camera preview.

% =====================================================================
\section{Inference}

\subsection{Whole frame, the default}

The image is letterboxed to \code{image\_size}, styled, and restored to its exact
original dimensions. Nothing is cropped and nothing is stretched. Verified
landmark round-trip: worst positional error $0.0017$, output size identical to
input, reconstruction error $0.02$--$0.03$ out of $255$.

\subsection{Tiled inference}

Optional, for native-resolution detail. Overlapping tiles are crossfaded with a
raised-cosine window.

\paragraph{Match the tile to the training crop.} This is not a preference. A model
trained on $256$ px crops has learned $256$ px-crop feature statistics; a
different tile size is a distribution shift. \code{Session.infer} warns on a
mismatch rather than blocking, because the mismatch is occasionally useful.

\begin{table}[h]
\centering
\begin{tabular}{@{}lrrrl@{}}
\toprule
Tile & Overlap ratio & Tiles & MP & Note \\
\midrule
256 & 25\% & 336 & 22.0 & matches a 256 px run \\
256 & 50\% & 713 & 46.7 & default; smoother \\
512 & 25\% & 88  & 23.1 & matches a 512 px run \\
1024& 25\% & 20  & 21.0 & matches a 1024 px run \\
\bottomrule
\end{tabular}
\caption{Tiled cost on a $4032\times3024$ source. What matters is the overlap
\emph{ratio}: the three $25\%$ rows all land near $22$ MP, and going to $50\%$
roughly doubles it.}
\label{tab:tiled}
\end{table}

\paragraph{Borders do not fade.} A Hann window tapers to zero at its edges, so a
naive crossfade fades the frame border towards black. Two mechanisms prevent it,
so no pre-padding is required: the window carries a positive floor of $10^{-3}$,
and edge tiles extend past the frame with the chosen padding, so the frame border
is never a tile border. Verified with a $255$-valued border strip, which survives
at $254.6/255$ with a worst-case error of $1.0$; a fade would show near zero.

The blend uses the same window weight in numerator and denominator. Clamping only
the denominator scales the output, which turned all four corners black in an
earlier version.

\subsection{Downscaling, and the training/inference gap}

\code{max\_edge} shrinks the long edge while preserving aspect ratio, before
tiling. It never upscales. Aspect ratio is preserved to within $0.005$ on the
sizes that matter.

\noindent
Degradation modes narrow the \emph{input} distribution to degraded night
photographs, which makes learning easier but can widen the gap to a real user
photograph. Three mitigations are in place, in order of importance: wide
per-image randomisation, the identity path, and the \code{strong} profile. An
explicit inference-time input normalisation was considered and \emph{rejected}: it
risks double-degrading a user image, and the training-side mitigations achieve
the same without inference cost or risk.

% =====================================================================
\section{Core ML boundary}

Export produces a \code{.mlpackage} or \code{.mlmodel}. All $28$ architecture
$\times$ head combinations have been verified to convert, with tensor-I/O mean
absolute error $0.00030$ against PyTorch.

Several implementation choices exist specifically because of this boundary, and
should not be ``cleaned up'':

\begin{itemize}[leftmargin=1.4em]
  \item \code{nn.MultiheadAttention} is unused. It lowers to a fused operator
        that fails conversion; attention uses explicit query/key/value matmuls.
  \item \code{Lut3DHead.table} is a flat $(s^3, 3)$ tensor. Per-channel gathers
        break Core ML scatter lowering.
  \item \code{grid\_sample} is unused anywhere. It returned near-centre values
        regardless of the grid in testing, and the 3D LUT uses explicit trilinear
        gathers.
  \item \code{HATBlock.self.grid} and the \code{EdgeAttention} window clamp use
        Python integers precomputed in \code{\_\_init\_\_}. Under
        \code{torch.jit.trace}, \code{x.shape[1]} becomes a zero-dimensional
        tensor, which breaks \code{max}, \code{min} and
        \code{adaptive\_avg\_pool2d}.
\end{itemize}

% =====================================================================
\section{Verification}
\label{sec:reproduce}

\subsection{Test suite}

$264$ tests pass. Beyond shape and export checks, the suite pins the
behaviours that would otherwise degrade silently:

\begin{itemize}[leftmargin=1.4em]
  \item the content-loss split, including that the $L_1$ half is insensitive to
        the styled image and the edge half is not;
  \item that the shared warp yields byte-identical source and target from an
        identical input;
  \item that validation degradation is invariant across epochs, and \emph{is}
        variant across them for training;
  \item that \code{sameset} never pairs an image with itself, and that a
        single-image pool falls back to \code{paired} with a warning rather than
        reaching \code{randrange(0)};
  \item that \code{best.pt} is not written in \code{sameset};
  \item an AST check that no module unpacks a generator call without
        \code{run\_generator}, which would break every newer architecture.
\end{itemize}

\subsection{Degradation probes}

Three scripts were used to derive the numbers in this report and are the basis
for the claims in \S\ref{sec:degrade}:

\begin{itemize}[leftmargin=1.4em]
  \item a stage-isolation probe, confirming each stage moves luminance in the
        intended direction;
  \item a halo probe, which is what falsified the claim that resampling removes
        halation;
  \item a real-photograph probe, reporting high-frequency energy, local contrast,
        tonal range, saturation and warmth per profile.
\end{itemize}

\subsection{Practical guidance}

\begin{description}[leftmargin=1.4em,style=nextline]
  \item[Target domain first.] A PatchGAN learns whatever distribution the target
        folder contains. If that folder is wrong, no architecture, schedule or
        tiling will help, and the failure is silent. Measure the target domain's
        mean luminance before spending GPU hours; for a Night Flash target it
        should sit below $0.20$, and daylit photography sits above $0.35$.
  \item[Watch the strip, not the loss.] The discriminator loss is not a quality
        signal.
  \item[Two arms, two GPUs.] At $4$ M parameters a single T4 is not saturated
        and data parallelism would add synchronisation cost without throughput.
        Two independent runs answer a real question that one run cannot.
  \item[No TPU.] The pipeline uses explicit gathers, \code{torch.jit} tracing for
        export, and a PatchGAN. A TPU would require an XLA rewrite of the
        augmentation path, and Core ML export would still happen on CPU.
\end{description}

\end{document}
